\documentclass[11pt,letterpaper]{article}

\usepackage[margin=0.75in]{geometry}
\usepackage{amsmath,amssymb}
\usepackage{booktabs,array,longtable,tabularx}
\usepackage{enumitem}
\usepackage{xcolor}
\usepackage{fancyhdr}
\usepackage{microtype}
\usepackage[hidelinks]{hyperref}

\definecolor{hkustblue}{RGB}{0,83,155}
\definecolor{lightblue}{RGB}{232,242,250}

\setlength{\parindent}{0pt}
\setlength{\parskip}{0.45em}
\setlength{\headheight}{14pt}
\setlist[itemize]{leftmargin=1.6em,itemsep=0.2em,topsep=0.2em}

\pagestyle{fancy}
\fancyhf{}
\fancyhead[L]{HKUST IEDA 1250}
\fancyhead[R]{Summer 2026 Final Exam Rubric}
\fancyfoot[C]{\thepage}

\newcommand{\E}{\mathbb{E}}
\newcommand{\rubricsection}[2]{%
  \section*{\color{hkustblue}Question #1 \normalfont\normalsize (#2 points)}
}
\newcolumntype{L}[1]{>{\raggedright\arraybackslash}p{#1}}
\newcolumntype{C}[1]{>{\centering\arraybackslash}p{#1}}

\begin{document}

\begin{center}
  {\LARGE\bfseries IEDA 1250, Summer 2026}\par
  \vspace{0.2em}
  {\Large Final Exam Marking Rubric}\par
  \vspace{0.4em}
  {\normalsize Total: 120 points}
\end{center}

\section*{General marking guidance}

\begin{itemize}
  \item Award full credit for any mathematically equivalent method with correct
  reasoning, even when notation differs from the official solution.
  \item Accept exact fractions or sufficiently accurate decimals.  Minor rounding
  differences should not be penalized.
  \item Apply follow-through credit: after one arithmetic error, grade later steps
  using the student's resulting value when the method remains correct.  Do not
  penalize the same error twice.
  \item Unsupported final answers receive only the points explicitly assigned to the
  final result; requested graphical or dynamic-programming work must be shown.
  \item For Questions 6 and 7, use the formula printed in the Summer 2026 paper,
  whose first denominator is $\lambda(\lambda+\mu_1)$.
  \item For Question 7, accept any valid stable numerical instance that proves the
  claimed strict comparison; the official instance is only one example.
\end{itemize}

\section*{Answer-key summary}

\begin{longtable}{L{0.10\textwidth}L{0.81\textwidth}}
\toprule
Question & Key result \\
\midrule
1 & Player 1: $(3/8,5/8)$; Player 2: $(1/2,1/2)$; game value to Player 1: $1/2$.\\
2 & $p^*=(0,1/3,2/3,0,0,0)$, $q^*=(1/2,1/2,0,0,0,0)$, $v=6$; the two stated LPs are primal and dual.\\
3 & Study-day allocation $(2,1,3,1)$; maximum grade points $23$.\\
4 & Location sequence $(A,B,C,B,B)$; maximum net profit $37$.\\
5 & $\pi=(16,8,4,2,1)/31$, utilization $15/31$, loss $1/31$, $\E[N]=26/31$, $\lambda_e=60/31$, $\E[T]=13/30$; capacity wins in (g), speed wins in (h).\\
6 & $\E[N]=1/[A(1-\rho)^2]$.\\
7 & One valid instance: $(\lambda,\mu_1,\mu_2)=(5,1,10)$, giving $\E[N]_{2}=1331/978>1=\E[N]_{\text{fast only}}$.\\
\bottomrule
\end{longtable}

\rubricsection{1}{20}

\begin{longtable}{L{0.12\textwidth}L{0.72\textwidth}C{0.08\textwidth}}
\toprule
Item & Expected work & Points \\
\midrule
1(a) & Defines $x=\Pr(\text{Player 1 uses }S_1)$ and obtains both lines
$L_1(x)=4x-1$ and $L_2(x)=2-4x$. & 3\\
1(a) & Correct graphical interpretation: Player 1 maximizes the lower envelope of
the two payoff lines. & 3\\
1(a) & Finds the intersection $x=3/8$, value $v=1/2$, and reports Player 1's
strategy $(3/8,5/8)$. & 3\\
1(a) & Makes Player 1 indifferent and obtains Player 2's strategy
$(1/2,1/2)$. & 3\\
1(b) & Constructs Player 2's correctly oriented payoff table
$\bigl(\begin{smallmatrix}-3&1\\2&-2\end{smallmatrix}\bigr)$. & 3\\
1(b) & Obtains $H_1(q)=2-5q$, $H_2(q)=3q-2$ and shows the graphical intersection
at $(q,v_2)=(1/2,-1/2)$. & 3\\
1(b) & Recovers $x=3/8$ and explicitly checks the sign relationship
$v_2=-v_1=-1/2$. & 2\\
\midrule
& \textbf{Question 1 total} & \textbf{20}\\
\bottomrule
\end{longtable}

\textit{Graphing note:} If the student gives correct algebra but no graph or envelope
interpretation, withhold up to 3 points across the graphical-work rows.

\rubricsection{2}{20}

\begin{longtable}{L{0.12\textwidth}L{0.72\textwidth}C{0.08\textwidth}}
\toprule
Item & Expected work & Points \\
\midrule
2(a) & Correctly removes dominated strategies: $A_5$ by $A_2$, $A_6$ by $A_3$;
$B_3,B_4$ by $B_1$ and $B_5,B_6$ by $B_2$. & 2\\
2(a) & Forms the reduced $4\times2$ game and obtains $q_1=q_2=1/2$ and $v=6$
from the relevant $A_2$--$A_3$ intersection. & 2\\
2(a) & Makes $B_1$ and $B_2$ indifferent and obtains
$p_2=1/3$, $p_3=2/3$. & 2\\
2(a) & Reports the full six-component strategies and checks that unused rows and
columns cannot improve either player's payoff. & 2\\
2(b) & Defines $p_i$ and $v_A$ and gives the objective $\max v_A$. & 1\\
2(b) & Writes all six full-table payoff constraints in the correct direction.  Award
$0.5$ per correct constraint, up to 3. & 3\\
2(b) & Includes $\sum_i p_i=1$, $p_i\ge0$, and $v_A$ unrestricted. & 1\\
2(c) & Rewrites the payoff constraints as
$v_A-\sum_i a_{ij}p_i\le0$ and assigns nonnegative dual variables $q_j$. & 2\\
2(c) & Introduces unrestricted $v_B$ for the primal equality and derives the dual
objective $\min v_B$. & 2\\
2(c) & Derives $\sum_j a_{ij}q_j\le v_B$ for all six rows and
$\sum_jq_j=1$ from the free variable $v_A$. & 2\\
2(c) & States that the resulting LP, including signs, is identical to the given
Player B LP. & 1\\
\midrule
& \textbf{Question 2 total} & \textbf{20}\\
\bottomrule
\end{longtable}

\rubricsection{3}{20}

\begin{longtable}{L{0.12\textwidth}L{0.72\textwidth}C{0.08\textwidth}}
\toprule
Item & Expected work & Points \\
\midrule
Setup & Defines a valid backward-DP value such as $G_i(s)$ and respects the
``at least one day per remaining course'' feasibility condition. & 2\\
Course 4 & Correct terminal-stage values $G_4(1,2,3,4)=(6,7,9,9)$. & 3\\
Course 3 & Shows the four maximizations and obtains
$G_3(2,3,4,5)=(8,10,13,14)$. & 5\\
Course 2 & Shows the four maximizations and obtains
$G_2(3,4,5,6)=(13,15,18,19)$. & 5\\
Course 1 & Evaluates
$G_1(7)=\max\{22,23,21,20\}=23$. & 3\\
Backtrack & Correctly backtracks to $(x_1,x_2,x_3,x_4)=(2,1,3,1)$ and verifies
$5+5+7+6=23$. & 2\\
\midrule
& \textbf{Question 3 total} & \textbf{20}\\
\bottomrule
\end{longtable}

\rubricsection{4}{20}

\begin{longtable}{L{0.12\textwidth}L{0.72\textwidth}C{0.08\textwidth}}
\toprule
Item & Expected work & Points \\
\midrule
Setup & Correctly interprets $(L,r)$, subtracts travel cost on a change, updates
$r$, and excludes choosing $L$ when $r=2$. & 2\\
Day 5 & Correct terminal continuation and Day 5 values:
$V_5(A,1)=7$, $V_5(A,2)=7$, $V_5(B,1)=10$, $V_5(B,2)=6$,
$V_5(C,1)=8$, $V_5(C,2)=8$, with valid best actions. & 3\\
Day 4 & Correct values $(16,14,13,13,15,15)$ for states
$(A,1),(A,2),(B,1),(B,2),(C,1),(C,2)$ and valid best actions. & 4\\
Day 3 & Correct values $(21,21,23,23,25,19)$ in the same state order and valid
best actions. & 4\\
Day 2 & Correct values $(29,29,32,30,30,30)$ in the same state order and valid
best actions. & 3\\
Day 1 and backtrack & Evaluates initial choices as $(37,34,32)$, selects $A$, and
backtracks to the unique optimal sequence $(A,B,C,B,B)$. & 2\\
Final value & Correctly shows daily net contributions $(8,6,8,5,10)$ and total
net profit $37$. & 2\\
\midrule
& \textbf{Question 4 total} & \textbf{20}\\
\bottomrule
\end{longtable}

\textit{Tie note:} At off-path states, accept either valid tied action, including
$B$ or $C$ at $V_5(C,1)$, $A$ or $B$ at $V_4(B,1)$, and $A$ or $C$ at
$V_3(A,1)$.

\rubricsection{5}{20}

\begin{longtable}{L{0.12\textwidth}L{0.72\textwidth}C{0.08\textwidth}}
\toprule
Item & Expected work & Points \\
\midrule
5(a) & Draws states $0,1,2,3,4$ with rightward rate $\lambda$ through state 4 and
leftward rate $\mu$ through state 0; recognizes arrivals at 4 are rejected. & 2\\
5(b) & Uses local/global balance to obtain $\pi_{n+1}=\rho\pi_n$ with
$\rho=1/2$. & 2\\
5(b) & Normalizes and reports
$(\pi_0,\ldots,\pi_4)=(16,8,4,2,1)/31$. & 2\\
5(c) & Utilization $1-\pi_0=15/31$. & 1\\
5(d) & States PASTA and obtains the loss fraction $\pi_4=1/31$.  One point for
PASTA/explanation and one for the value. & 2\\
5(e) & Computes $\E[N]=\sum n\pi_n=26/31$. & 2\\
5(f) & Computes the effective arrival rate
$\lambda_e=\lambda(1-\pi_4)=60/31$. & 1\\
5(f) & Applies Little's Law to admitted jobs and obtains
$\E[T]=(26/31)/(60/31)=13/30$. & 2\\
5(g) & Correct loss probabilities $1/511$ after capacity doubling and $1/341$
after speed doubling; concludes capacity doubling is better. & 2\\
5(h) & Correct loss probabilities $6561/242461$ after capacity doubling and
$81/6505$ after speed doubling; concludes speed doubling is better. & 2\\
5(i) & Explains that added space is especially effective at light load, whereas at
heavier load service capacity addresses the congestion bottleneck directly. & 2\\
\midrule
& \textbf{Question 5 total} & \textbf{20}\\
\bottomrule
\end{longtable}

\rubricsection{6}{10}

\begin{longtable}{L{0.12\textwidth}L{0.72\textwidth}C{0.08\textwidth}}
\toprule
Item & Expected work & Points \\
\midrule
Law & States and correctly applies Little's Law $\E[N]=\lambda\E[T]$. & 3\\
Substitution & Substitutes the provided $\E[T]$ and cancels $\lambda$. & 2\\
Result & Obtains $\E[N]=1/[A(1-\rho)^2]$. & 3\\
Definitions & Retains the printed definitions of $A$ and
$\rho=\lambda/(\mu_1+\mu_2)<1$ without swapping $\mu_1$ and $\mu_2$. & 2\\
\midrule
& \textbf{Question 6 total} & \textbf{10}\\
\bottomrule
\end{longtable}

\rubricsection{7}{10}

\begin{longtable}{L{0.12\textwidth}L{0.72\textwidth}C{0.08\textwidth}}
\toprule
Item & Expected work & Points \\
\midrule
Instance & Gives parameters with $\mu_2>\mu_1$, $\lambda<\mu_1+\mu_2$, and
$\lambda<\mu_2$.  Official example: $(5,1,10)$. & 2\\
Two servers & Correctly computes $\rho$, $A$, and two-server $\E[N]$.  For the
official example: $\rho=5/11$, $A=163/66$, $\E[N]=1331/978$. & 3\\
Fast only & Uses the $M/M/1$ formula $\E[N]=\lambda/(\mu_2-\lambda)$.  The
official example gives $1$. & 2\\
Comparison & Shows a strict improvement after disconnection; for the official
example, $1<1331/978$. & 1\\
Intuition & Explains how fast-server-first routing can send a job to a very slow
server even when waiting briefly for the fast server would be better, while the fast
server alone still has sufficient spare capacity. & 2\\
\midrule
& \textbf{Question 7 total} & \textbf{10}\\
\bottomrule
\end{longtable}

\section*{Point-total check}

\[
20+20+20+20+20+10+10=\boxed{120}.
\]

\end{document}
